%
% Documento: Metodologia
%

\chapter{Metodologia}
\label{cap:metodologia}

\section{Caracterização da pesquisa}

Esta pesquisa é aplicada, quantitativa e experimental. Ela reúne o desenvolvimento da arquitetura e do firmware de um protótipo RFID UHF e a execução de uma campanha de medições com o conjunto formado por leitor, antena e tag. A unidade de análise é a combinação entre leitor, antena, etiqueta, material de fixação, geometria e condição de teste.

O inventário portátil constitui o cenário principal. O ponto fixo é tratado como detecção de presença em uma zona de leitura; a direção do movimento não é inferida com uma única antena. Os percentuais apresentados neste trabalho resultam dos dados coletados durante a campanha experimental.

\section{Configuração experimental}

A configuração utilizada nos ensaios é formada pelo YPD-R200, pela antena Airplux APCA8090 e pela etiqueta Avery Dennison Smartrac AD-239 M730. A \autoref{tab:materiais-prototipo} reúne as especificações de referência e as condições elétricas registradas.

\begin{table}[htbp]
    \centering
    \caption{Componentes empregados na campanha experimental.}
    \label{tab:materiais-prototipo}
    \small
    \begin{tabularx}{\textwidth}{p{3.2cm}X}
        \toprule
        \textbf{Componente} & \textbf{Especificação de referência} \\
        \midrule
        Leitor & YPD-R200, variante com potência máxima nominal de 26~dBm, aproximadamente 398~mW no conector de RF, EPC Gen2/ISO 18000-6C, alimentação de 3,6 a 5,5~V, UART TTL a 115200~bit/s e conexão SMA \cite{yanpodo2021r200}. \\
        Antena & Airplux APCA8090, ganho nominal de 5~dBi, polarização linear, frequências centrais informadas de 915~MHz e 866,5~MHz, impedância de 50~$\Omega$ e elemento cerâmico de 80~$\times$~80~$\times$~7~mm \cite{airplux2026apca8090}. \\
        Tag & Avery Dennison Smartrac AD-239 M730, chip Impinj M730, de 860 a 960~MHz, ISO/IEC 18000-63, EPC Class 1 Gen2, memória EPC de 128 bits e aplicação recomendada em superfícies não metálicas \cite{avery2022ad239m730}. \\
        Controlador & Arduino Uno, com uso da abstração \texttt{Stream} e de \texttt{SoftwareSerial} na captura dos quadros. \\
        Alimentação e lógica & Foram utilizados, em momentos distintos, uma fonte externa de 5~V/1~A e o fornecimento de energia pela conexão USB do Arduino. O caminho Arduino TX $\rightarrow$ R200 RXD requer adequação de 5~V para 3,3~V; o emprego efetivo desse adequador não ficou confirmado de modo inequívoco. \\
        \bottomrule
    \end{tabularx}
    \fonte{Elaborado pelo autor a partir da montagem experimental e dos documentos dos fabricantes (2026).}
\end{table}

O leitor, a antena e o Arduino foram montados em bancada. A campanha principal ocorreu em 13 de julho de 2026, no Laboratório de Práticas Digitais, ambiente fechado de aproximadamente 50~m$^2$. Durante essa campanha, havia cerca de cinco pessoas e quatro objetos metálicos próximos. Essa configuração é compatível com uma sala institucional de uso cotidiano e foi adotada como cenário realista de inventário patrimonial, e não como substituta de uma bancada destinada à caracterização eletromagnética. Utilizaram-se UART a 115200~bit/s e o comando nominal de potência máxima de 26~dBm, correspondente a aproximadamente 398~mW no conector do módulo. Não houve instrumento ou comando de leitura posterior que confirmasse a potência efetivamente aplicada ou irradiada. Em momentos distintos da coleta principal foram empregados a fonte externa de 5~V/1~A e o arranjo alimentado pela conexão USB do Arduino. Após as medições, verificou-se que a alternância entre as formas de alimentação não produziu alteração relevante nos resultados obtidos. A altura da antena e da tag foi mantida em 150~cm nas medições de distância, orientação e material. Foram aguardados 3~s antes da primeira tentativa e 2~s entre tentativas.

A região, o canal e o comportamento de salto de frequência não foram lidos de volta pelo firmware durante a campanha. Assim, embora a potência nominal, a antena e a faixa regulamentar aplicável sejam conhecidas, os dados coletados não demonstram por si só a subfaixa efetivamente utilizada nem a potência irradiada. Essas ausências foram mantidas como ameaças à validade, em vez de se presumir conformidade a partir de uma configuração nominal.

\section{Montagem e coleta}

Nos ensaios de distância, orientação e material, o leitor e a antena permaneceram fixos e buscou-se manter as demais condições enquanto se variava o fator de cada bloco. As condições foram executadas em grupos e não houve aleatorização formal da ordem. No inventário em sala, o conjunto foi deslocado pelo operador ao longo de um percurso curto. O registro serial foi capturado junto com o EPC observado, o tempo até a primeira leitura e observações sobre leituras externas, \textit{timeout} e quadros inválidos. O RSSI fornecido pelo módulo foi registrado em formato bruto, sem conversão para dBm, porque a campanha não incluiu calibração instrumental suficiente para isso. A \autoref{fig:arquitetura-prototipo} sintetiza o fluxo entre a etiqueta, o leitor, o Arduino e os dados coletados.

\begin{figure}[htbp]
    \centering
    \caption{Arquitetura funcional usada na campanha experimental.}
    \label{fig:arquitetura-prototipo}
    \setlength{\tabcolsep}{3pt}
    \begin{tabular}{ccccc}
        \fbox{\shortstack{AD-239 M730\\EPC Gen2}}
        & $\rightleftarrows$ &
        \fbox{\shortstack{Antena\\5 dBi}}
        & $\rightleftarrows$ &
        \fbox{\shortstack{YPD-R200\\RF + protocolo}}
        \\
        &&&& $\downarrow$ UART TTL \\
        &&&& \fbox{\shortstack{Arduino\\aquisição}} $\longrightarrow$
        \fbox{\shortstack{Dados\\coletados}}
    \end{tabular}
    \fonte{Elaborado pelo autor (2026).}
\end{figure}

O firmware carregado no Arduino foi desenvolvido para esta campanha. Durante uma janela de inventário, a implementação compara EPCs pelo conteúdo e pelo comprimento, cria uma entrada apenas para cada EPC ainda não visto, incrementa a contagem das repetições e separa tags cadastradas de leituras externas. Portanto, a deduplicação usada no inventário foi executada no Arduino; ela não foi inferida posteriormente apenas a partir da quantidade total de quadros.

\subsection{Ensaios de distância e orientação}

Os ensaios abrangeram as distâncias horizontais de 0,5~m, 1,0~m, 1,5~m, 2,0~m e 5,0~m e as orientações relativas de 0$^\circ$, 45$^\circ$ e 90$^\circ$. O ângulo de 0$^\circ$ correspondeu ao alinhamento de referência entre as orientações lineares da antena e da etiqueta. Em 0,5~m foi executado apenas esse alinhamento; nas demais distâncias foram executados os três ângulos. Como as condições reuniram etapas de coleta com números diferentes de repetições, cada combinação contém de 15 a 45 tentativas, totalizando 375 observações na análise conjunta.

No subconjunto com registros individuais, as tags TAG1, TAG2 e TAG3 foram empregadas em conjuntos fixos com papelão, madeira e plástico, respectivamente. Essa associação impede separar o efeito do exemplar do efeito de seu suporte. Parte das demais observações foi preservada como contagem de detecções por combinação. A janela de medição foi de 3~s, com antena e tag na mesma altura de 150~cm, e considerou-se bem-sucedida a tentativa na qual o EPC esperado foi identificado pelo menos uma vez. Nos gráficos, as taxas são estratificadas por distância e ângulo. As sínteses marginais são apenas descritivas, pois o número de repetições, a composição das etiquetas e a disponibilidade de registros individuais não são iguais em todas as condições.

\subsection{Ensaios de material}

O ensaio de material fixou a distância em 1,5~m e utilizou TAG1 em três condições: papelão, metal direto e metal com espaçador de 10~mm. Cada condição foi repetida 5 vezes. A janela de 3~s foi mantida e o espaçador foi utilizado para verificar se a separação física recuperava a leitura em condição crítica.

\subsection{Inventário em sala}

O inventário utilizou cinco tags presentes simultaneamente e cinco rodadas de observação. Cada janela teve 25~s, tempo superior ao dos demais ensaios porque o objetivo era acumular EPCs distintos e não apenas validar uma leitura isolada. O operador repetiu um percurso curto padronizado. A métrica diagnóstica foi a cobertura por rodada, enquanto o critério operacional foi identificar as cinco tags na mesma passagem. Essa população constituía a etapa inicial de um escalonamento gradual até aproximadamente 50 etiquetas. A ampliação dependia de cobertura integral na etapa inicial; como nenhuma rodada foi completa, o número de tags não foi aumentado. A disposição usual do laboratório foi preservada para representar a operação cotidiana em uma sala institucional de aproximadamente 50~m$^2$. Como suas dimensões lineares, o trajeto e as coordenadas das tags não foram registrados com precisão, esse bloco permite avaliar o desempenho entre as cinco rodadas, mas não reproduzir exatamente a distribuição espacial da campanha. A \autoref{tab:plano-ensaios} resume os fatores, as repetições e as saídas efetivamente analisadas.

\begin{table}[htbp]
    \centering
    \caption{Cenários e parâmetros executados na campanha experimental.}
    \label{tab:plano-ensaios}
    \footnotesize
    \begin{tabularx}{\textwidth}{p{3.2cm}p{4.2cm}p{3.0cm}X}
        \toprule
        \textbf{Cenário} & \textbf{Fatores e níveis} & \textbf{Repetições} & \textbf{Saída analisada} \\
        \midrule
        Distância e orientação & 0,5 a 5,0~m e 0$^\circ$, 45$^\circ$ e 90$^\circ$, conforme as combinações executadas & 15 a 45 por combinação & Detecções, taxa de leitura e IC 95\% de Wilson; tempo e RSSI apenas no subconjunto com registros individuais. \\
        Material & Papelão, metal direto e metal com espaçador de 10~mm & 5 por condição & Taxa de leitura e RSSI bruto. \\
        Inventário & Cinco tags em sala, janela de 25~s & 5 rodadas & EPCs distintos, cobertura, omissões e duplicatas. \\
        \bottomrule
    \end{tabularx}
    \fonte{Elaborado pelo autor a partir do protocolo experimental e dos dados coletados (2026).}
\end{table}

\FloatBarrier

Não foram executados ensaios independentes de passagem em porta, variação de altura, variação controlada de potência ou comparação entre antenas. Consequentemente, esses fatores são tratados apenas como requisitos ou possibilidades de continuidade e não originam resultados neste trabalho.

\section{Métricas e tratamento dos dados coletados}

Para $N_t$ tentativas e $N_v$ detecções, a taxa de sucesso é:

\begin{equation}
    T_s = \frac{N_v}{N_t}\times 100\%.
    \label{eq:taxa-sucesso}
\end{equation}

Em uma rodada com $N_p$ tags presentes e $N_d$ EPCs distintos detectados, a cobertura é:

\begin{equation}
    C = \frac{N_d}{N_p}\times 100\%.
    \label{eq:cobertura}
\end{equation}

Para $N_r$ rodadas de inventário e $N_c$ rodadas nas quais todos os EPCs esperados foram identificados, a taxa de inventários completos é:

\begin{equation}
    T_c = \frac{N_c}{N_r}\times 100\%.
    \label{eq:taxa-inventarios-completos}
\end{equation}

A taxa de não detecção e a taxa de omissão são os complementos dessas métricas:

\begin{equation}
    T_{nl} = \frac{N_t-N_v}{N_t}\times 100\% = 100\%-T_s,
    \label{eq:taxa-nao-leitura}
\end{equation}

\begin{equation}
    T_o = \frac{N_p-N_d}{N_p}\times 100\% = 100\%-C.
    \label{eq:taxa-omissao}
\end{equation}

Para $N_q$ eventos válidos e $N_d$ EPCs distintos, a taxa de duplicação é:

\begin{equation}
    T_{dup} = \frac{N_q-N_d}{N_q}\times 100\%.
    \label{eq:duplicacao}
\end{equation}

Os indicadores foram calculados a partir dos dados coletados durante os ensaios. Para cada proporção de detecção e para a proporção de rodadas completas foi calculado o intervalo de confiança de 95\% pelo método de Wilson, adequado inclusive quando a taxa observada é 0\% ou 100\%. As tabelas marginais de distância e orientação reúnem todas as observações disponíveis em cada nível, mas não representam efeitos isolados, porque a distribuição de ângulos, etiquetas, suportes e repetições é desbalanceada. A interpretação principal, portanto, considera as séries estratificadas dos gráficos. Para tempo até a primeira leitura e RSSI bruto são apresentados média e desvio-padrão amostral somente entre as detecções válidas com registros individuais. No inventário, a cobertura média é acompanhada do desvio-padrão entre rodadas e de um intervalo exploratório de 95\% para a média com distribuição $t$ de Student. O intervalo da cobertura média descreve a fração média de etiquetas identificadas, não a probabilidade de terminar o inventário. Para a decisão operacional, utiliza-se $T_c$, pois uma cobertura parcial não atende ao requisito de conferir todos os bens em uma passagem. Esses intervalos descrevem a incerteza da amostra, mas devem ser lidos como exploratórios: as janelas reutilizam as mesmas etiquetas e o mesmo ambiente, de modo que não constituem observações plenamente independentes. Não foram aplicados testes de significância, pois o número reduzido de repetições e essa estrutura de dependência não sustentariam inferências amplas. Como o RSSI foi armazenado em formato bruto, ele é tratado apenas como indicador relativo do conjunto formado por leitor, tag e ambiente.

\section{Validade prática e limitações}

A principal limitação é o tamanho da amostra e a realização dos ensaios em uma única sala. A escolha de um laboratório institucional de uso cotidiano, com ocupação e objetos metálicos próximos, fortalece a validade prática da avaliação: essas condições são compatíveis com ambientes nos quais ocorre a conferência patrimonial e integram o funcionamento esperado, não sendo tratadas como erros experimentais. Como os elementos do ambiente não foram posicionados e variados independentemente, os resultados expressam o desempenho conjunto no cenário observado, sem permitir quantificar separadamente as contribuições do multipercurso e da movimentação de pessoas. Nos ensaios de distância e orientação, a identidade da etiqueta e o material de suporte permaneceram associados, o que também impede atribuir diferenças entre TAG1, TAG2 e TAG3 apenas ao exemplar. Parte dessas medições foi preservada apenas como contagens agregadas, sem RSSI, tempo individual, ordem das tentativas ou diagnóstico serial por tentativa. Também não houve instrumentação de RF para calibrar o RSSI ou caracterizar o canal.

A montagem elétrica e a configuração de RF também possuem lacunas de rastreabilidade. O emprego do adequador de nível entre D11 e RXD não ficou confirmado de modo inequívoco, e região, canal, salto de frequência e potência efetiva não foram verificados por leitura posterior ou instrumentação. Além disso, quadros inválidos e \textit{timeouts} foram observados com \texttt{SoftwareSerial} a 115200~bit/s, mas não houve analisador lógico para separar perda serial de ausência de resposta no enlace de RF.

Outra limitação é o uso de uma única antena linear. Isso foi suficiente para caracterizar a referência do protótipo, mas não permite comparar diretamente outras geometrias. Polarização circular ou diversidade de polarização trata mais diretamente a sensibilidade angular observada. Uma antena omnidirecional é uma hipótese distinta para ampliar cobertura azimutal na varredura portátil e pode também aumentar multipercurso ou leituras externas. Ambas as possibilidades exigem comparação sob o mesmo protocolo.

Registros fotográficos da montagem foram preservados como documentação auxiliar. Como não constituem medição quantitativa e a arquitetura está representada esquematicamente, sua inclusão no corpo da monografia não foi necessária para interpretar os resultados.

Por fim, o firmware ainda não é um sistema patrimonial completo. Ele deduplica os EPCs durante a janela e sustenta a campanha atual, mas uma versão operacional futura deve incluir persistência de dados, filtros mais robustos e validação em mais de uma sessão e ambiente de coleta.
